\documentclass[10pt,a4paper]{article}
\usepackage[utf8]{inputenc}
\usepackage[T1]{fontenc}
\usepackage[english]{babel}
\usepackage{geometry}
\geometry{margin=1.8cm}
\usepackage{enumitem}
\usepackage{booktabs}
\usepackage{fancyhdr}
\usepackage{lastpage}
\usepackage{tcolorbox}

\pagestyle{fancy}
\fancyhf{}
\rfoot{\thepage/\pageref{LastPage}}
\lhead{\small TOEFL ITP 2026 --- Day 4}
\rhead{\small MPrometeo}
\renewcommand{\headrulewidth}{0.4pt}
\setlength{\parindent}{0pt}
\setlength{\parskip}{5pt}

\begin{document}

\begin{center}
{\LARGE \textbf{DAY 4 --- Monday Mar 2}}\\[4pt]
{\large Reading Strategy}\\[2pt]
{\normalsize Based on Pynchon's \emph{Slow Learner}
--- Introduction (real text)}\\[2pt]
\rule{\textwidth}{0.8pt}
\end{center}

% ============ CONTEXT ============
\section*{Source Text}

\begin{tcolorbox}[colback=white, colframe=black]
All exercises below are derived from the actual Introduction to
\emph{Slow Learner} (1984) by Thomas Pynchon --- the autobiographical
essay before the stories, NOT the stories themselves. Direct quotes
are used throughout.
\end{tcolorbox}

% ============ STRATEGIES ============
\section*{Reading Strategies for TOEFL ITP}

\subsection*{Strategy 1: Skimming (30 seconds per passage)}

Read ONLY the first and last sentence of each paragraph. This gives
you the ``map.''

\textbf{Applied to Pynchon's Introduction:}

His first paragraph opens with:
\begin{quote}\small
``...written between 1958 and 1964. Four of them I wrote when I was
in college.''
\end{quote}
And closes with:
\begin{quote}\small
``...how comfortable would I feel about lending him money, or for
that matter even stepping down the street to have a beer and talk
over old times?''
\end{quote}

From just these two sentences you know: he's reflecting on old
work and feels distant from his younger self. That's the map ---
everything between is detail.

\subsection*{Strategy 2: Scanning (for specific info)}

When a question says ``According to the passage...'' do NOT reread
everything:
\begin{enumerate}[nosep]
\item Find the \textbf{keyword} in the question
\item \textbf{Scan} for that word or its synonym
\item Read only \textbf{2 sentences around it}
\item Choose your answer
\end{enumerate}

\subsection*{Strategy 3: Vocabulary in Context}

Pynchon uses words in unexpected ways. TOEFL tests this constantly.

\textbf{Method:} Replace the word with each option. Which one makes
the sentence mean the same thing?

Example from his text: ``these stories will still be of use with
all their \textbf{flaws} intact.''

(A) strengths (B) defects (C) pages (D) themes $\to$ Answer: (B)

\subsection*{Strategy 4: Timing}

\begin{center}
\begin{tabular}{ll}
\toprule
\textbf{TOEFL ITP Reading} & \textbf{Your target} \\
\midrule
50 questions, 5 passages & 55 min total \\
= 11 min per passage & 4 min read + 7 min questions \\
Stuck $>$30 seconds? & Mark it, move on, come back \\
\bottomrule
\end{tabular}
\end{center}

% ============ PASSAGE 1 ============
\section*{Reading Passage 1: Pynchon on His Own Failures ($\sim$220 words)}

\begin{tcolorbox}[colback=white, colframe=black]
\small
Pynchon opens the Introduction to \emph{Slow Learner} with a
confession that rereading his early stories produced physical
discomfort: ``My first reaction, rereading these stories, was oh my
God, accompanied by physical symptoms we shouldn't dwell upon.'' His
second impulse was ``some kind of a wall-to-wall rewrite,'' but
eventually he settled into what he calls ``middle-aged tranquility''
--- a pretense of clarity about who he was as a young writer.

He warns readers directly: ``there are some mighty tiresome passages
here, juvenile and delinquent too.'' Yet his hope is that the
stories, ``pretentious, goofy and ill-considered as they get,'' will
still serve as ``illustrative of typical problems in entry-level
fiction, and cautionary about some practices which younger writers
might prefer to avoid.''

About his first published story, ``The Small Rain,'' Pynchon admits
he failed to trust his own material. Instead of letting the main
character's problem generate the story naturally, he felt compelled
to add ``a whole extra overlay of rain images and references to
\emph{The Waste Land} and \emph{A Farewell to Arms}.'' His guiding
principle at the time, he confesses, was a motto he invented himself:
``Make it literary'' --- which he calls ``a piece of bad advice I
made up all by myself and then took.''
\end{tcolorbox}

\subsection*{Questions (time yourself: 7 min)}

\begin{enumerate}[label=\textbf{\arabic*.}, leftmargin=*]

\item What was Pynchon's first reaction to rereading his stories?
\begin{enumerate}[label=(\Alph*)]
\item Pride and satisfaction
\item Physical discomfort and dismay
\item Indifference
\item Excitement to publish them again
\end{enumerate}

\item ``Pretentious'' in paragraph 2 is closest in meaning to:
\begin{enumerate}[label=(\Alph*)]
\item humble \item trying to appear more important than reality
\item funny \item brief
\end{enumerate}

\item Why does Pynchon criticize ``The Small Rain''?
\begin{enumerate}[label=(\Alph*)]
\item It was too short
\item He added unnecessary literary references instead of trusting
the story
\item The characters were too realistic
\item It was never published
\end{enumerate}

\item ``Cautionary'' most likely means:
\begin{enumerate}[label=(\Alph*)]
\item encouraging \item serving as a warning
\item entertaining \item confusing
\end{enumerate}

\item What was Pynchon's self-invented motto?
\begin{enumerate}[label=(\Alph*)]
\item ``Write what you know''
\item ``Make it literary''
\item ``Keep it simple''
\item ``Trust the reader''
\end{enumerate}
\end{enumerate}

\textbf{Answers:} 1-B, 2-B, 3-B, 4-B, 5-B

% ============ PASSAGE 2 ============
\section*{Reading Passage 2: On Faking Knowledge ($\sim$230 words)}

\begin{tcolorbox}[colback=white, colframe=black]
\small
One of Pynchon's most striking confessions in the Introduction
concerns his relationship with vocabulary and research. He admits to
a ``specific piece of wrong procedure'': browsing through the
thesaurus to find words that ``sounded cool, hip, or likely to
produce an effect, usually that of making me look good, without then
taking the trouble to go and find out in the dictionary what they
meant.'' He calls this behavior ``stupid'' and mentions it only
hoping others might ``profit from my error.''

This extends beyond vocabulary to factual knowledge. Pynchon notes
that young writers often ``think we know everything --- or to put it
more usefully, we are often unaware of the scope and structure of
our ignorance. Ignorance is not just a blank space on a person's
mental map. It has contours and coherence, and for all I know rules
of operation as well.''

He illustrates this with an embarrassing example from ``Entropy'':
he used the French phrase \emph{grippe espagnole}, thinking it
referred to some kind of post-World War I spiritual malaise. It
actually means ``Spanish influenza'' --- the literal flu epidemic.
His lesson: ``corroborate one's data, in particular those acquired
casually, such as through hearsay or off the backs of record albums.''
The parallel to TOEFL preparation is direct: never use a word you
cannot define precisely.
\end{tcolorbox}

\subsection*{Questions}

\begin{enumerate}[label=\textbf{\arabic*.}, leftmargin=*]

\item What ``wrong procedure'' does Pynchon confess to?
\begin{enumerate}[label=(\Alph*)]
\item Plagiarizing other authors
\item Using words from a thesaurus without checking their meaning
\item Writing too slowly
\item Avoiding all research
\end{enumerate}

\item According to Pynchon, ignorance:
\begin{enumerate}[label=(\Alph*)]
\item Is simply an empty space
\item Has structure, contours, and coherence
\item Only affects young people
\item Cannot be overcome
\end{enumerate}

\item What was Pynchon's mistake with \emph{grippe espagnole}?
\begin{enumerate}[label=(\Alph*)]
\item He spelled it wrong
\item He thought it meant spiritual malaise when it meant Spanish flu
\item He used it in the wrong story
\item He invented the phrase himself
\end{enumerate}

\item ``Corroborate'' is closest in meaning to:
\begin{enumerate}[label=(\Alph*)]
\item invent \item ignore \item verify/confirm \item translate
\end{enumerate}

\item What lesson does Pynchon draw for writers?
\begin{enumerate}[label=(\Alph*)]
\item Never do research
\item Always verify your information, especially casually acquired facts
\item Use as many French phrases as possible
\item Trust everything you read on album covers
\end{enumerate}
\end{enumerate}

\textbf{Answers:} 1-B, 2-B, 3-B, 4-C, 5-B

% ============ PASSAGE 3 ============
\section*{Reading Passage 3: On the Beat Generation and Literary
Loyalty ($\sim$220 words)}

\begin{tcolorbox}[colback=white, colframe=black]
\small
Pynchon describes his college generation as caught between two
literary worlds. On one side stood the established modernist
tradition they were studying in class. On the other, the Beat
writers --- Kerouac, Ginsberg, and others --- offered a radical
alternative. ``We were at a transition point, a strange post-Beat
passage of cultural time, with our loyalties divided,'' he writes.

The appeal of the Beats was powerful. Pynchon recalls discovering
the first issue of the \emph{Evergreen Review} in a Norfolk,
Virginia bookstore in 1956 while serving in the Navy: ``It was an
eye-opener.'' He describes fellow sailors who ``would sit in circles
on the deck and sing perfectly, in parts, all those early rock 'n'
roll songs, who played bongos and saxophones, who had felt honest
grief when Bird and later Clifford Brown died.''

Yet Pynchon acknowledges that his generation were ``onlookers: the
parade had gone by and we were already getting everything
secondhand.'' They adopted Beat postures and props, but the original
energy had already passed. Still, he insists, what remained was ``a
sane and decent affirmation of what we all want to believe about
American values.'' And he identifies Kerouac's \emph{On the Road}
as ``one of the great American novels'' --- a judgment he states
without qualification or irony.
\end{tcolorbox}

\subsection*{Questions}

\begin{enumerate}[label=\textbf{\arabic*.}, leftmargin=*]

\item Pynchon's generation was ``caught between'':
\begin{enumerate}[label=(\Alph*)]
\item Science and art
\item Modernist tradition and Beat literature
\item Europe and America
\item Music and writing
\end{enumerate}

\item ``Eye-opener'' most likely means:
\begin{enumerate}[label=(\Alph*)]
\item A surgical procedure
\item Something that was surprisingly revealing
\item A type of magazine
\item An alarm clock
\end{enumerate}

\item Pynchon says his generation was:
\begin{enumerate}[label=(\Alph*)]
\item The leaders of the Beat movement
\item Onlookers who received the movement secondhand
\item Completely unaware of the Beats
\item Opposed to all forms of rebellion
\end{enumerate}

\item What does Pynchon call \emph{On the Road}?
\begin{enumerate}[label=(\Alph*)]
\item An overrated book
\item One of the great American novels
\item A travel guide
\item A failed experiment
\end{enumerate}

\item It can be inferred that Pynchon views the Beat
movement with:
\begin{enumerate}[label=(\Alph*)]
\item Pure contempt
\item Affection and respect despite distance
\item Total indifference
\item Anger and resentment
\end{enumerate}
\end{enumerate}

\textbf{Answers:} 1-B, 2-B, 3-B, 4-B, 5-B

% ============ SHADOWING ============
\section*{Shadowing Exercise (15 min --- Gemini Live)}

\begin{tcolorbox}[colback=white, colframe=black,
    title=Tell Gemini Live:]
``Read this aloud. I repeat sentence by sentence.
Score my pronunciation and fluency.''
\end{tcolorbox}

\begin{quote}\small
``My first reaction, rereading these stories, was oh my God,
accompanied by physical symptoms we shouldn't dwell upon. It is only
fair to warn even the most kindly disposed of readers that there are
some mighty tiresome passages here, juvenile and delinquent too. My
specific piece of wrong procedure was to browse through the thesaurus
and note words that sounded cool, without taking the trouble to find
out what they meant.''
\end{quote}

\textbf{Pronunciation targets:}
\begin{itemize}[nosep]
\item \textbf{accompanied} --- a-CUM-pa-need (stress on 2nd)
\item \textbf{juvenile} --- JOO-ve-nile (stress on 1st)
\item \textbf{thesaurus} --- the-SOR-us (stress on 2nd)
\item \textbf{procedure} --- pro-SEE-jur (stress on 2nd)
\item \textbf{illustrative} --- i-LUS-tra-tiv (stress on 2nd)
\end{itemize}

% ============ PYNCHON'S ADVICE ============
\section*{Bonus: Pynchon's Advice Applied to YOUR TOEFL}

\begin{tcolorbox}[colback=white, colframe=black,
    title=What Pynchon Teaches You About English]
\begin{enumerate}[nosep]
\item \textbf{Don't fake vocabulary.} Pynchon: ``browse through the
thesaurus and note words that sounded cool, without taking the
trouble to find out what they meant.'' On TOEFL, use only words you
truly understand.

\item \textbf{Know the shape of your ignorance.} ``Ignorance is not
just a blank space. It has contours and coherence.'' Identify your
weak grammar patterns --- don't just ``study more.''

\item \textbf{Verify your data.} ``Corroborate one's data, in
particular those acquired casually.'' On TOEFL reading, always go
back to the text --- never answer from memory alone.

\item \textbf{Be a slow learner.} The title says it all. 28 days
of honest, consistent work beats pretending to know more than you do.
\end{enumerate}
\end{tcolorbox}

\end{document}